% ============================================================
% paper/sections/08_phase_logic.tex
% Phase Logic: SUCCESS, INERTIA, COLLAPSE, STOP
% ============================================================

\section{Phase Logic}

\subsection{Purpose of Phase Classification}

Phase logic in $K_{\mathrm{EA}}$ provides a coarse-grained diagnostic
classification of admissible enterprise states. Phases do not represent
stages of maturity or progress. They are labels assigned based on the
evaluation of threshold components and continuity conditions.

The purpose of phase logic is to determine which kinds of diagnostic
statements remain admissible.

\subsection{Precedence Rules}

Phases are evaluated according to a strict precedence order defined in
\texttt{ETS.K\_EA.v1.1}:

\[
\text{STOP} \; \succ \; \text{COLLAPSE} \; \succ \; \text{INERTIA} \;
\succ \; \text{SUCCESS}.
\]

Higher-precedence phases override lower-precedence classifications.
This prevents contradictory diagnoses.

\subsection{STOP Phase}

The STOP phase represents termination of diagnosis. It is entered when
any of the following conditions hold:

\begin{itemize}
  \item $f_{\text{exist}} > 0$,
  \item $f_{\text{stop}} > 0$,
  \item observability required to evaluate thresholds is lost.
\end{itemize}

In the STOP phase, further diagnosis is undefined within the model. No
claims regarding recovery, improvement, or evolution are admissible.

\subsection{COLLAPSE Phase}

The COLLAPSE phase indicates irreversible loss of structural coherence
while diagnosis remains momentarily defined.

Formally, COLLAPSE holds when:
\begin{itemize}
  \item $f_{\text{collapse}} > 0$,
  \item STOP conditions are not met.
\end{itemize}

COLLAPSE does not imply immediate non-existence, but it indicates that
the enterprise continuum can no longer satisfy stability conditions
over the internal time scale.

\subsection{INERTIA Phase}

The INERTIA phase characterizes states in which the enterprise remains
structurally admissible but cannot transition between admissible states
without external intervention.

Formally, INERTIA holds when:
\begin{itemize}
  \item $f_{\text{inertia}} > 0$,
  \item STOP and COLLAPSE conditions are not met.
\end{itemize}

Diagnosis remains defined, but internal dynamics are constrained.

\subsection{SUCCESS Phase}

The SUCCESS phase is defined by the absence of threshold violations:

\[
f_k \leq 0 \quad \text{for all components } k.
\]

SUCCESS does not imply optimality, efficiency, or desirability. It
indicates only that the enterprise continuum satisfies admissibility
conditions for diagnosis.

\subsection{Allowed Transitions}

The ETS specifies the following allowed transitions:

\begin{itemize}
  \item SUCCESS $\rightarrow$ INERTIA, COLLAPSE, STOP,
  \item INERTIA $\rightarrow$ SUCCESS via external $\Psi$, COLLAPSE,
        STOP,
  \item COLLAPSE $\rightarrow$ STOP,
  \item STOP $\rightarrow$ (none).
\end{itemize}

Transitions are constrained by threshold evaluations and parent
compatibility conditions.

\subsection{Non-Teleological Interpretation}

Phase logic does not encode a preferred direction of movement. There is
no assumption that SUCCESS should be maintained or restored.

Phases are diagnostic labels, not objectives.
